% Generated draft: requires mathematical and source review.
\documentclass[11pt]{article}
\usepackage[margin=25mm]{geometry}
\usepackage{fontspec}
\setmainfont{Latin Modern Roman}
\usepackage{amsmath,amssymb,mathtools,unicode-math}
\setmathfont{Latin Modern Math}
\usepackage{enumitem,needspace,xurl,hyperref}
\setlength{\parindent}{0pt}
\setlength{\parskip}{5pt}
\setlength{\emergencystretch}{3em}
\newcommand{\N}{\mathbb{N}}
\newcommand{\Z}{\mathbb{Z}}
\newcommand{\Q}{\mathbb{Q}}
\newcommand{\R}{\mathbb{R}}
\newcommand{\C}{\mathbb{C}}
\newcommand{\F}{\mathbb{F}}
\newcommand{\A}{\mathbb{A}}
\newcommand{\PP}{\mathbb{P}}
\newcommand{\E}{\mathbb{E}}
\newcommand{\eps}{\varepsilon}
\newcommand{\dd}{\,\mathrm{d}}
\begin{document}
\section*{UnsolvedMath ID: 30006462}
\textbf{Source problem code:} catalogue entry
\subsection{Definitions and mathematical statement}
Let \(\mathbb H=\{x+iy:x\in\R,\ y>0\}\), with hyperbolic measure \(\dd\mu=y^{-2}\dd x\,\dd y\) and nonnegative hyperbolic Laplacian
\[
\Delta=-y^2\left(\frac{\partial^2}{\partial x^2}+\frac{\partial^2}{\partial y^2}\right).
\]
A congruence subgroup \(\Gamma\leq\mathrm{SL}_2(\Z)\) is a subgroup containing
\[
\Gamma(m)=\{\gamma\in\mathrm{SL}_2(\Z):\gamma\equiv I\pmod m\}
\]
for some integer \(m\geq1\). It acts on \(\mathbb H\) by fractional linear transformations. In the classical scalar-valued, weight-zero interpretation, a Maass cusp form for \(\Gamma\) is a nonzero smooth function \(f:\mathbb H\to\C\) that is \(\Gamma\)-invariant, belongs to \(L^2(\Gamma\backslash\mathbb H,\dd\mu)\), has zero constant Fourier coefficient at every cusp, and satisfies \(\Delta f=\lambda_f f\).

The supplied target is Selberg's eigenvalue conjecture: for every congruence subgroup \(\Gamma\leq\mathrm{SL}_2(\Z)\) and every such Maass cusp form with nonzero eigenvalue,
\[
\lambda_f\geq\frac14.
\]
Equivalently, there is no such form with \(0<\lambda_f<\tfrac14\); these are the exceptional Maass cusp forms relevant here.

The scalar-valued, weight-zero convention and the Laplacian normalization have been made explicit to interpret the stated threshold. The supplied record does not specify a weight, character, or multiplier system, so extensions to those settings are not asserted. Its original excerpt specifically discusses \(\Gamma_0(q)=\{\left(\begin{smallmatrix}a&b\\c&d\end{smallmatrix}\right)\in\mathrm{SL}_2(\Z):q\mid c\}\), for positive integers \(q\), whereas its clean statement quantifies over congruence subgroups generally. The excerpt presents exceptional-spectrum losses and on-average large sieve substitutes as motivation; it does not supply a sufficiently legible complete inequality to formulate a separate large sieve problem. The conjecture's open status is supplied dated curation, not an independently verified conclusion.
\subsection{Short English statement}
Selberg's eigenvalue conjecture predicts that every nonzero Laplace eigenvalue of a classical weight-zero Maass cusp form for any congruence subgroup is at least \(1/4\). Thus no exceptional cuspidal eigenvalue lies in \((0,1/4)\). The supplied excerpt motivates this through large sieve estimates for \(\Gamma_0(q)\); its damaged formulas and unspecified automorphy conventions require human review.
\subsection{Sources}
\begin{itemize}
\item Local registry: \texttt{data/unsolvedmath/problems.json}, permanent UnsolvedMath ID \(30006462\), title ``Selberg Eigenvalue Conjecture for Congruence Groups,'' source problem code \texttt{catalogue entry}, in the Oberwolfach Reports Open Problems collection. The supplied record attributes the conjecture to Selberg; its proposer field is empty and its proposed year is \(2025\).
\item Supplied source citation: \emph{Analytic Number Theory (2026)}, Oberwolfach Reports, DOI \texttt{10.4171/owr/2025/51}, pp.\ 2701--2764. Source URL and supplied literature reference ``Oberwolfach report 51/2025'': \url{https://doi.org/10.4171/owr/2025/51}. The citation's year \(2026\) and the report designation \(51/2025\) are retained as supplied; their relationship has not been verified.
\item Supplied literature reference: ``Pascadi, exceptional Maass forms,'' \url{https://doi.org/10.1017/fmp.2026.10025}. No fuller bibliographic details are supplied.
\item Literature assessment checked at \texttt{2026-08-22}: the supplied curation labels the problem open, says Selberg's eigenvalue conjecture remains unproved, and says current unconditional bounds are strictly weaker than \(1/4\). This is supplied dated curation derived from the OWR material, not independently verified literature assessment. The record explicitly requests independent verification of time-sensitive status and cited claims.
\item The supplied statement-status label is \texttt{corrected\_verified}. Its accompanying verification note says that only clear formula delimiters, notation, grammar, and typography were repaired, with mathematical content and scope preserved. This label does not establish independent mathematical or source verification of the present draft.
\end{itemize}
\end{document}
